% Source: https://artofproblemsolving.com/community/c3h1675626_let_a_a1a2__a100_and_b__12__50_determine_the_number_of_surjec
% Stable identity: aops-topic-1675626-q01
A nondecreasing function is determined by how many elements go to each value: if $n_j$ is the number of $a_i$ with $f(a_i)=j$, then the values must be $1$ repeated $n_1$ times, then $2$ repeated $n_2$ times, and so on, so $f$ is determined by $(n_1,\dots,n_{50})$. Conversely every $50$-tuple of nonnegative integers with $n_1+\cdots+n_{50}=100$ gives such a function. So we count these tuples (stars and bars).

\textbf{(a)} Surjective means every $n_j\ge1$. Writing $100$ as an ordered sum of $50$ positive integers is the same as choosing $49$ of the $99$ gaps between $100$ stars in a row to place bars, so there are
\[
\boxed{\dbinom{99}{49}}
\]
such functions.

\textbf{(b)} Now $n_j\ge0$; with $m_j=n_j+1\ge1$ the sum becomes $150$, so the number is
\[
\boxed{\dbinom{149}{49}} .
\]
